\section{本章小结}
\label{sec:opt-fit-summary}

章首的两条曲线在观测点上几乎一样，在点与点之间却截然不同。本章的回答是：读数只确定函数的一部分，其余部分由我们写进目标的偏好、所选的表示和所用的算法共同决定，每一种决定都需要说明条件。

观测首先划出可见与不可见的方向。有限个有界线性读数只看见表示元张成的有限维子空间，余项完全不可见；读数能否区分全部候选，由候选差集与观测零空间是否只交于零决定。正则化在不可见的方向上作出选择，却不增加读数的识别能力；残差代价决定怎样对待偏离的读数，硬约束与有限正则参数是不同的问题，目标中的二次惩罚与更新规则中的解耦衰减在一般的预条件下也求解不同的问题。

表示把无限维的选择化成有限维的计算。固定基函数时，二次拟合是增广最小二乘，唯一性由读数与惩罚共同的零方向决定；稳健代价借切线上界化为一串加权最小二乘；非线性参数借残差或输出的线性化求步，再由原目标检验接受。候选取遍Hilbert空间时，表示定理把搜索限制到表示元的张成空间，Gram矩阵给出读数与范数；整体函数值观测则在核积分算子的谱坐标中逐方向配方，梯度观测化为到闭子空间的投影。在这些二次问题中，函数存在且唯一，系数却可以沿Gram矩阵的零空间不唯一。

误差要按来源与范数分开。数据扰动经过谱乘子传到最优函数，最坏的放大发生在特征值与$n\lambda$相当的方向；正则参数趋于零时，有限观测下的解趋于最小范数最小二乘解，整体观测下的解在$L^2$中收敛，而在$H$中只在数据的正谱部分属于$H$时收敛；表示误差与优化误差在同一目标值中可以相加，再借目标值差与函数距离的关系换成距离。最后，当目标有许多最小点或没有有限最小点时，初值、预条件、参数化与算法本身还会作出选择。

\begin{table}[!htbp]
  \centering
  \small
  \begin{tabular}{@{}>{\raggedright\arraybackslash}p{22mm}>{\raggedright\arraybackslash}p{23mm}>{\raggedright\arraybackslash}p{32mm}>{\raggedright\arraybackslash}p{15mm}@{}}
    \toprule
    观测与候选 & 求解工具 & 存在与唯一 & 本章位置\\
    \midrule
    有限读数，固定基 & 增广最小二乘、IRLS & $\ker\bm X\cap\ker\bm R=\{\bm0\}$时唯一 & 第~\ref{sec:opt-fit-linear}节\\
    有限读数，非线性参数 & GN、LM、广义GN & 局部；比值或回溯检验下降 & 第~\ref{sec:opt-fit-nonlinear}节\\
    有限读数，Hilbert空间 & 表示定理、Gram矩阵 & $\lambda>0$时函数唯一，系数可不唯一 & 第~\ref{sec:opt-fit-hilbert}节\\
    整体函数值 & 谱坐标逐方向配方 & $\lambda>0$时唯一 & 第~\ref{sec:opt-fit-function-values}小节\\
    梯度场 & 闭子空间投影 & 固定边界唯一；自由边界唯一到常数 & 第~\ref{sec:opt-fit-gradient-observations}小节\\
    无正则的多解问题 & 梯度下降的轨迹 & 极限由初值与预条件决定 & 第~\ref{sec:opt-fit-implicit-bias}节\\
    \bottomrule
  \end{tabular}
  \caption{本章各类拟合问题的求解工具与结论。“唯一”均指最优函数；系数表示是否唯一另行判断。}
  \label{tab:opt-fit-summary}
\end{table}

函数唯一、系数唯一与统计意义上的可识别，是三种不同的判断：前两者是本章对确定性问题的回答，后者涉及数据的生成机制，留给统计推断与学习理论。本章能够完整求解，靠的是目标为二次、可以配方和投影；一般的非二次泛函，例如含梯度的非线性积分目标，就需要直接扰动函数、计算变化率并分析积分目标的最优条件与存在性，这是下一章的任务。
